\section{Síntesis y ampliación}

Esta clase une la historia de los modelos de lenguaje, el aprendizaje no supervisado, la generación multimodal y la code evaluation en una sola metodología:

\begin{enumerate}
  \item \textbf{La calidad de la generación tiene niveles}: el léxico, la sintaxis, la semántica, los hechos y el éxito en la tarea no pueden reducirse a «si se parece o no»;
  \item \textbf{La perplexity solo mide la predicción de la distribución}: no demuestra por sí sola la veracidad, la controlabilidad ni la downstream utility;
  \item \textbf{El autoregressive objective aporta autosupervisión densa}: el analysis by synthesis puede producir una representation transferible, pero también puede aprender shortcuts;
  \item \textbf{GPT-1/2/3 cambiaron la interfaz de las tareas}: el fine-tuning, la zero-shot continuation y las in-context demonstrations ampliaron sucesivamente la reutilización;
  \item \textbf{Lo multimodal exige redefinir el token y la metric}: iGPT y DALL-E demuestran que la recipe es transferible, no que la estructura del dominio carezca de importancia;
  \item \textbf{El código necesita un functional verifier}: el overlap de cadenas no puede sustituir a la semántica de ejecución;
  \item \textbf{Pass@k mide el rare success dentro de la distribución}: pass@k y pass@1 corresponden a distintos candidate budgets y a distintas experiencias de producto;
  \item \textbf{La temperature debe calibrarse según k}: una temperatura baja busca la calidad de una sola muestra; una temperatura alta, con el respaldo de un verifier, puede ampliar la cobertura de la búsqueda;
  \item \textbf{El sampling solo puede amplificar la masa de probabilidad existente}: si la probabilidad del programa correcto es cero, ninguna cantidad de samples sirve;
  \item \textbf{La production reliability requiere un proceso completo}: el sandbox, los tests, el static analysis, la review, la security y el maintenance no pueden sustituirse con un benchmark.
\end{enumerate}

\begin{importantbox}{Tarea de Generative-System Evaluation}
Elige una tarea de texto, imagen o código: define la model distribution, la decoding policy, el candidate budget, el verifier y el final selector; reporta a la vez pass@1/quality, pass@k/diversity, cost, latency y la failure taxonomy. Para la tarea de código, deriva el estimador de pass@k y diseña tres tipos de tests ocultos; para la tarea de texto o imagen, explica por qué tu verifier no puede ser un oracle perfecto.
\end{importantbox}

\begin{knowledgebox}{Autoevaluación final: hay que distinguir tres «existencias»}
\textbf{Existir en la distribución}: alguna muestra correcta tiene probabilidad distinta de cero; \textbf{poder encontrarse}: la temperature, el valor de $k$ y la search actuales permiten muestrearla; \textbf{poder reconocerse}: los tests o el ranker pueden seleccionarla de forma confiable. La mayoría de las diferencias entre productos de modelos generativos ocurren entre estos tres niveles.
\end{knowledgebox}

\subsection{Lecturas adicionales}

{\small
\begin{itemize}[nosep,leftmargin=*]
  \item \href{https://www.cs.utoronto.ca/~ilya/pubs/2011/LANG-RNN.pdf}{\textbf{Generating Text with Recurrent Neural Networks}}: RNN language generation.
  \item \href{https://arxiv.org/abs/1602.02410}{\textbf{Exploring the Limits of Language Modeling}}: large LSTM language modeling.
  \item \href{https://arxiv.org/abs/1704.01444}{\textbf{Learning to Generate Reviews and Discovering Sentiment}}: sentiment neuron y unsupervised representation.
  \item \href{https://cdn.openai.com/research-covers/language-unsupervised/language_understanding_paper.pdf}{\textbf{GPT-1}} y \href{https://cdn.openai.com/better-language-models/language_models_are_unsupervised_multitask_learners.pdf}{\textbf{GPT-2}}: de pretrain/fine-tune a zero-shot.
  \item \href{https://arxiv.org/abs/2005.14165}{\textbf{GPT-3}}: few-shot/in-context evaluation y human detection study.
  \item \href{https://arxiv.org/abs/2006.10738}{\textbf{iGPT}}: generative pretraining from pixels.
  \item \href{https://arxiv.org/abs/2102.12092}{\textbf{DALL-E}}: zero-shot text-to-image generation.
  \item \href{https://arxiv.org/abs/2107.03374}{\textbf{Codex}}: HumanEval, pass@k y code generation evaluation.
\end{itemize}
}

\end{document}
